% Agent context and source locators: ../../../internal/section_notes/02_goals.md
\section{Research Questions, Aims, and Staged Goals}\label{sec:research-questions-aims-and-staged-goals}

\subsection{Framework contribution}\label{subsec:framework-contribution}
\begin{itemize}[leftmargin=*]
\item Develop and assess a CFD--experiment order-reduction framework that produces a fast, interpretable, physics-based reduced-order model (ROM).
\item Use a custom physical 1D solver as the principal model; CFD and experiments inform reduction, closure selection, parameter learning, and assessment.
\item Proposed aims: (1) reduction and development-informed closures; (2) parameter learning and observation-consistent temperature reconstruction; (3) independent assessment and applicability.
\item Organizing question: what information and reduced physics are needed to reproduce observed temperatures and then predict new conditions with defensible uncertainty? Keep the development-case demonstration distinct from the independent predictive test.
\item For each aim, define a target output, permitted inputs, comparator, and assessment evidence.

\item \topiclabel{closest-prior-methods-and-completed-groundwork}\textbf{Closest prior methods and completed groundwork.} The closest work includes thermal-network structure/\allowbreak{}identification and physical interpretation, goal-oriented reduction/\allowbreak{}calibration/\allowbreak{}design, embedded discrepancy, conjugate/\allowbreak{}stochastic loop dynamics, conservative/\allowbreak{}stable learned reduction and similarity.\cite{fierro2026cfdoneD,raue2009identifiability}
\begin{itemize}
\item Compare the proposed approach with existing thermal-network, goal-oriented, discrepancy, conjugate and stochastic, and structure-preserving methods, as well as completed preliminary modeling studies.
\end{itemize}
\item \topiclabel{candidate-contribution-packages-and-rejection-tests}\textbf{Candidate contributions.} Evaluate supported-target and anchoring analysis, finite-time retained-physics adequacy, qualified transfer, and the thermal consequences of local flow as possible contributions within the reduction framework.
\begin{itemize}
\item For each candidate, specify the added-value claim, closest prior method, fair baseline, information requirements, failure test, and fallback. The dissertation need not pursue every branch.
\end{itemize}
\item \topiclabel{which-additional-result-could-advance-the-framework}\textbf{Potential methodological advances.} Candidate outcomes include a reproducible development-informed reduction and closure rule, observation-consistent temperature reconstruction with an assessed domain, experiment--CFD parameter learning under shared uncertainty, or a qualified physical limit.
\begin{itemize}
\item Test a precise closure/\allowbreak{}reconstruction/\allowbreak{}learning procedure or qualified applicability limit; use temperature/\allowbreak{}circulation/\allowbreak{}duty relevance rather than the novelty of software integration.
\end{itemize}
\item \topiclabel{method-procedure-physics-and-facility-demonstration}\textbf{Method, procedure, physics and facility demonstration.} Method development supplies a new or meaningfully improved rule and shows why/\allowbreak{}when it works.
\begin{itemize}
\item State which kind of contribution each work package would establish and what evidence separates it from implementation or a retrospective fit.
\end{itemize}
\item \topiclabel{a-useful-negative-result}\textbf{A useful negative result.} State a model or evidence class, a consequential ambiguity or adequacy boundary, and reproducible evidence that advances beyond established scale or rank limitations and prior preliminary findings.
\begin{itemize}
\item A failed fit or unavailable CFD alone is insufficient; a negative result should identify a consequential limit on prediction or physical interpretation.
\end{itemize}
\end{itemize}

\subsection{Staged research goals}\label{subsec:staged-research-goals}

\begin{itemize}[leftmargin=*]
\item Stage 1: steady heat removal/\allowbreak{}losses, circulation, temperature extrema and supported operating margins.
\item Stage 2: transitions between operating conditions, thermal and hydraulic response, settling, and time-dependent margins.
\item Stage 3: operational assistance and operator-facing model/\allowbreak{}data dashboard; collaborate with Ben Booth and NeutronOS. Controls and unattended operation need a separate agreed scope.
\item Uncertainty estimation and propagation are central to all three stages.
\item The immediate feasibility demonstration emphasizes experimental-data interpretation and temperature reproduction, building on the completed master\textquotesingle{}s work.
\end{itemize}

\subsection{Spatial questions and boundaries}\label{subsec:spatial-questions-and-boundaries}
\begin{itemize}[leftmargin=*]
\item Temperature first: centerline and radial profiles, with angular dependence where supported; richer velocity mapping is a later capability.
\item Overall circulation and local recirculation are separate outputs; investigate local extent and thermal consequence with qualified CFD, without inventing local experimental measurements.
\end{itemize}
